\textbf{Identification almost surely.} Fix \(a\in\{0,1\}\). By \(\mathrm{ass:conditional-exchangeability}\),
\[
\mathbb E_P[Y^a\mid X,A=a]=\mathbb E_P[Y^a\mid X]
\]
where the left-hand conditional mean is defined. By \(\mathrm{ass:strict-overlap}\), the event \(A=a\) has positive conditional probability for \(P_X\)-almost every covariate value, so the displayed identity holds \(P_X\)-almost surely. By \(\mathrm{ass:consistency}\), \(Y=Y^a\) on \(\{A=a\}\), and therefore
\[
\mu_a(X)=\mathbb E_P[Y\mid X,A=a]=\mathbb E_P[Y^a\mid X]
\]
almost surely. Subtracting the identities for \(a=0\) and \(a=1\), then invoking \(\mathrm{def:cate}\), gives
\[
\tau_P(X)=\mu_1(X)-\mu_0(X)=\mathbb E_P[Y^1-Y^0\mid X]
\]
almost surely.

\textbf{Uniqueness on the support.} Let \(g_1\) and \(g_2\) be continuous representatives of the same \(P_X\)-almost-sure function. If \(g_1(x_0)\ne g_2(x_0)\) at some \(x_0\in\operatorname{supp}(P_X)\), continuity gives a relatively open neighborhood \(O\) of \(x_0\) and \(\varepsilon>0\) such that \(|g_1-g_2|>\varepsilon\) on \(O\). The definition of support gives \(P_X(O)>0\), contradicting almost-sure equality. Hence the continuous representative is unique on \(\operatorname{supp}(P_X)\).

Two observed-law versions of \(\tau_P\) that agree \(P_X\)-almost surely have upper sets, as defined in \(\mathrm{def:upper-set}\), that differ only on a \(P_X\)-null set. Likewise, replacing either set argument by a set that differs on a \(P_X\)-null set leaves
\[
d_H(H,G_P;P)=\int_{H\triangle G_P}|\tau_P(x)-\theta|f_P(x)\,dx
\]
unchanged, because this is integration with respect to \(P_X\). This proves the observed-law interpretation and the asserted null-set invariance.